Cellular responses to drug treatments and genetic perturbations provide an experimental basis for comparing drug mechanisms and gene functions. The Connectivity Map linked perturbations with similar effects through transcriptional signatures~\cite{lamb2006}, and L1000 expanded the scale of expression profiling~\cite{subramanian2017}. Cell Painting records changes in cellular structure and morphology through multiplexed imaging, providing information complementary to expression profiles~\cite{bray2016,way2022}. These measurements allow researchers to search for functionally related drugs or genetic perturbations from their cellular responses. However, the same biological process can produce different responses across measurement modalities and cellular backgrounds, while experimental batches also affect similarities between profiles. Functional retrieval therefore requires identifying features of perturbation responses that connect measurements across these conditions.

CellProfiler extracts morphological features~\cite{stirling2021}, and DeepProfiler learns image representations~\cite{moshkov2024}. Multimodal methods incorporate additional information: CLOOME pairs images with chemical structures~\cite{sanchez2023}, while Semi-Clipped and TIDE transfer information between morphology and expression through cross-modal distillation~\cite{bendidi2025,chen2026}. Pairing images and expression profiles establishes correspondence between measurements of the same drug. Biological text and annotations additionally provide functional information about perturbations, as used by CellCLIP, MorphoCLIP and annotation-informed contrastive learning~\cite{lu2025,ilyosbekov2026,pwesombo2025}. Genetic imaging and matched chemical--genetic screens provide links between gene perturbations and drug responses~\cite{rohban2017,chandrasekaran2024}, allowing gene-function findings and chemical interventions to inform each other~\cite{li2026unipert}. For retrieval across these settings, representations should retain same-drug correspondence and connect different drugs and genetic perturbations by function.

To combine measurement pairing with functional associations, we present BioMoAX, which learns cellular perturbation representations from Cell Painting images, L1000 expression profiles, and chemical--genetic image pairs through five complementary contrastive objectives. In a separate compound projection space, the compound-paired objective aligns images and expression profiles of the same drug to learn cross-modal identity correspondence. In a shared functional projection space, image MOA and expression MOA objectives connect drugs with shared mechanisms within each modality, while a cross-modal MOA objective connects mechanism-related drugs across modalities. A gene--drug paired objective uses CRISPR and associated drug image pairs to incorporate genetic perturbations into this functional space. By separating compound identity from perturbation function and jointly modeling within-modality, cross-modal, and chemical--genetic relationships, BioMoAX supports retrieval of functionally related responses across measurement modalities and perturbation types (Fig.~1A).

We evaluated compound identity and within-cell-line MOA retrieval before extending queries across experimental batches, cell lines and perturbation types (Fig.~1B). In bidirectional U2OS--A549 retrieval, image queries identified mechanism-related image and expression profiles from the other cell line. HT29 images excluded from image training queried the U2OS/A549 library, achieving image-to-image (I2I) and image-to-expression (I2E) Top-5 MOA hit rates of 21.78\% and 18.69\%, compared with 8.45\% for CellProfiler and 5.35\% for ridge-mapped CLOOME, the strongest corresponding baselines. The I2E queries crossed both cellular background and measurement modality, using existing transcriptomic profiles as references for image-only queries. CRISPR queries for held-out genes achieved a Top-10 hit rate of 19.77\%, compared with 13.79\% for CellProfiler. KDR and DDR2 perturbations co-retrieved linifanib, while ALDH2 and AKR1B1 perturbations co-retrieved disulfiram. In both cases, the retrieved drug connects a reference-annotated gene--drug pair with an additional association supported by independent literature~\cite{karaman2008,yoshino2020npl4}.
